% Источник: лекция 8, с. 5--8; лекция 9, с. 9;
% лекция 10, с. 15--18 рукописного PDF.

\subsection{Счётные операции над множествами}

\begin{center}
    \begin{tikzpicture}[>=Stealth, every node/.style={font=\small}]
        \node (s) at (0,0) {Полукольцо};
        \node (r) at (3,0) {Кольцо};
        \node (sr) at (6.1,0.35) {$\sigma$-кольцо};
        \node (dr) at (6.1,-0.55) {$\delta$-кольцо};
        \draw[->] (s) -- (r);
        \draw[->] (r) -- (4.6,0) |- (sr);
        \draw[->] (4.6,0) |- (dr);
        \draw[->] (sr) -- (dr);
        \node (sa) at (0,-1.65) {Полуалгебра};
        \node (a) at (3,-1.65) {Алгебра};
        \node (sigma) at (6.1,-1.65) {$\sigma$-алгебра};
        \node (delta) at (6.1,-2.65) {$\delta$-алгебра};
        \draw[->] (sa) -- (a);
        \draw[->] (a) -- (sigma);
        \node at (6.1,-2.15) {$\Updownarrow$};
    \end{tikzpicture}
\end{center}

\begin{definition}[$\sigma$-алгебра с единицей $\Omega$]
    Система $\mathcal{F} \subset 2^\Omega$ называется
    \textbf{$\sigma$-алгеброй с единицей $\Omega$}, если
    \begin{enumerate}[label=\arabic*)]
        \item $\Omega \in \mathcal{F}$;
        \item $\mathcal{F}$ замкнута относительно операций
        $\cap$, $\cup$, $\setminus$, $\triangle$:
        \[
            A,B \in \mathcal{F}
            \implies A \star B \in \mathcal{F},
            \qquad \star \in \{\cap,\cup,\setminus,\triangle\};
        \]
        \item для любой последовательности $A_i \in \mathcal{F}$
        \[
            \bigcup_{i=1}^{\infty} A_i \in \mathcal{F}.
        \]
    \end{enumerate}
    При этом
    \[
        \varnothing = \Omega \setminus \Omega \in \mathcal{F},
        \qquad
        \bigcap_{i=1}^{\infty} B_i
        = \overline{\bigcup_{i=1}^{\infty}\overline{B_i}}
        \in \mathcal{F}
        \quad (B_i \in \mathcal{F}).
    \]
\end{definition}

\subsection{Наименьшая \texorpdfstring{$\sigma$}{σ}-алгебра}

\begin{definition}
    Пусть $\chi$ --- система множеств,
    $\Omega = \bigcup_{A\in\chi} A$.
    Наименьшая $\sigma$-алгебра $\sigma(\chi)$ определяется условиями:
    \begin{enumerate}[label=\arabic*)]
        \item $\chi \subset \sigma(\chi)$;
        \item $\sigma(\chi)$ --- $\sigma$-алгебра с единицей $\Omega$;
        \item для любой $\sigma$-алгебры $\mathcal{A}$
        с той же единицей $\Omega$
        \[
            \chi \subset \mathcal{A}
            \implies \sigma(\chi) \subset \mathcal{A}.
        \]
    \end{enumerate}
\end{definition}

\textbf{Построение:}
\[
    \sigma(\chi) = \bigcap_{\alpha} \mathcal{A}_\alpha,
\]
где $\mathcal{A}_\alpha$ --- все $\sigma$-алгебры
с единицей $\Omega$, содержащие $\chi$.

\subsection{Принцип подходящих множеств}

\begin{task}
    Для любого $A\in R(\chi)$ существует конечное покрытие
    элементами $\chi$:
    \[
        \exists A_1,\ldots,A_n\in\chi\colon
        A\subset\bigcup_{i=1}^{n} A_i.
    \]
\end{task}

\begin{solution}
    Рассмотрим
    \[
        R_0 = \left\{A\in R(\chi)\;\middle|\;
        \exists A_1,\ldots,A_n\in\chi\colon
        A\subset\bigcup_{i=1}^{n} A_i\right\}.
    \]
    Тогда
    \[
        R_0\subset R(\chi),\qquad
        \chi\subset R_0,\qquad R_0\text{ --- кольцо}.
    \]
    По минимальности $R(\chi)$:
    \[
        R(\chi)\subset R_0
        \quad\Longrightarrow\quad R_0=R(\chi).
    \]
\end{solution}

\[
    S\longrightarrow R(S)\longrightarrow\sigma(S),
    \qquad
    R(S)=\left\{\bigsqcup_{i=1}^{n}A_i\;\middle|\;
    n\in\N,\ A_i\in S\right\}.
\]

\subsection{Борелевская \texorpdfstring{$\sigma$}{σ}-алгебра}

\begin{definition}
    \textbf{Борелевской $\sigma$-алгеброй} на пространстве $A$
    называется наименьшая $\sigma$-алгебра, содержащая все
    открытые множества этого пространства.
    Обозначение: $\mathcal{B}(A)$.
\end{definition}

\subsubsection{Открытые множества на \texorpdfstring{$\mathbb{R}$}{R}}

\begin{theorem}
    Любое открытое множество $A\subset\R$ представляется
    в виде конечного или счётного дизъюнктного объединения
    открытых интервалов:
    \[
        A=\bigsqcup_{i\in I}(a_i;b_i),
        \qquad a_i,b_i\in\overline{\R},
    \]
    где $I$ --- не более чем счётное множество индексов.
\end{theorem}

\begin{Proof}
    Введём на $A$ отношение эквивалентности:
    \[
        x\sim y \quad\Longleftrightarrow\quad [x;y]\subset A.
    \]
    Здесь $[x;y]$ --- отрезок между $x$ и $y$.
    \begin{enumerate}[label=\arabic*)]
        \item Рефлексивность: $x\sim x$ для всех $x\in A$.
        \item Симметричность: $x\sim y\implies y\sim x$.
        \item Транзитивность:
        $x\sim y$, $y\sim z\implies x\sim z$.
    \end{enumerate}
    Разобьём $A$ на классы эквивалентности:
    \[
        A=\bigsqcup_{\alpha} A_\alpha.
    \]
    \textbf{1. $A_\alpha$ --- открытый интервал.}
    По определению отношения $A_\alpha$ является промежутком.
    Пусть, например, $A_\alpha=(a;b]$.
    Тогда $b\in A_\alpha\subset A$, поэтому
    существует $U_\eps(b)\subset A$.
    Пусть $(b;b+\eps)\subset A_\beta$.
    При достаточно малом $\eps>0$:
    \[
        b-\frac{\eps}{2}\sim b+\frac{\eps}{2},
    \]
    что противоречит тому, что $b$ --- правый конец $A_\alpha$.

    \textbf{2. Классов не более чем счётное число.}
    В каждом $(a_\alpha;b_\alpha)$ выберем рациональное число:
    \[
        q_\alpha\in\Q\cap(a_\alpha;b_\alpha).
    \]
    Интервалы не пересекаются, поэтому выбранные числа различны.
\end{Proof}

\subsubsection{Порождающие системы}

\begin{theorem}
    \[
        \begin{aligned}
            \mathcal{B}(\R)
            &=\sigma\bigl(\{(a;b)\}_{a,b\in\R}\bigr)
              =\mathcal{B}_2\\
            &=\sigma\bigl(\{[a;b]\}_{a,b\in\R}\bigr)
              =\mathcal{B}_3\\
            &=\sigma\bigl(\{(-\infty;x]\}_{x\in\R}\bigr)
              =\mathcal{B}_4\\
            &=\sigma\bigl(\{(a;b)\}_{a,b\in\Q}\bigr)
              =\mathcal{B}_5.
        \end{aligned}
    \]
\end{theorem}

\begin{Proof}
    \textbf{$\mathcal{B}(\R)=\mathcal{B}_4$.}
    Так как $(x;+\infty)$ открыто, то
    \[
        (-\infty;x]=\R\setminus(x;+\infty)\in\mathcal{B}(\R).
    \]
    Следовательно, $\mathcal{B}_4\subset\mathcal{B}(\R)$.

    Обратно, $(y;x]\in\mathcal{B}_4$ для всех $x,y\in\R$, и
    \[
        (a_i;b_i)=\bigcup_{n=1}^{\infty}
        \left(a_i;b_i-\frac1n\right]\in\mathcal{B}_4.
    \]
    Любое открытое множество --- конечное или счётное
    объединение интервалов. Поэтому
    $\mathcal{B}(\R)\subset\mathcal{B}_4$.

    \textbf{$\mathcal{B}_2=\mathcal{B}_5$.}
    Очевидно, $\mathcal{B}_5\subset\mathcal{B}_2$.
    Для любого $(a;b)$ выберем
    \[
        a_n,b_n\in\Q,\qquad a_n\downarrow a,
        \qquad b_n\uparrow b.
    \]
    Тогда
    \[
        (a;b)=\bigcup_{n=1}^{\infty}(a_n;b_n)\in\mathcal{B}_5.
    \]
    Следовательно, $\mathcal{B}_2\subset\mathcal{B}_5$.
\end{Proof}

\begingroup
\tcbset{myboxstyle/.append style={breakable=false}}
\begin{example}
    \[
        S=\bigl\{(a;b],\ (a;+\infty),\ (-\infty;b],\ \R
        \;\bigm|\; a,b\in\R\bigr\}.
    \]
    Для этой системы $\sigma(S)=\mathcal{B}(\R)$.
\end{example}
\endgroup
